\bookchapter{Capítulo 0: Olá, antes de começar}

Antes de explorarmos oficialmente, vamos falar sobre algumas regras básicas para garantir que compreende tudo o que se segue. Sem pressa, afiar a lâmina não atrasa o corte.

\booksection{I. O que este livro pode ajudar a fazer? (E o que não pode)}

Pode ajudá-lo a compreender sistematicamente como usar o celular, a entender a lógica por trás das funções e, quer seja utilizador Android ou Apple, a maioria das abordagens é comum. Vai acompanhá-lo em exercícios práticos, para que passe de “saber usar” a “compreender como usar”.

Não é o manual de um celular específico, nem notícias de tecnologia sobre os últimos lançamentos, e muito menos um guia para o tornar dependente do celular. Sinceramente, se espera “como assistir a vídeos com mais prazer no celular”, pode ficar desiludido — preocupamo-nos mais em como fazer com que o celular trabalhe verdadeiramente para si, e não o contrário.

\booksection{II. Explicação dos nossos “códigos”}

Para sermos mais claros, o livro utilizará algumas pequenas marcações, que funcionam como códigos de entendimento entre nós:

\begin{itemize}
  \item \textbf{[Nome do botão]}: representa um botão específico no tela ou o nome de uma App, como [Definições], [OK].
  \item \textbf{>}: representa a ordem das operações, por exemplo "tocar em [Definições] > [Som]" significa primeiro tocar em Definições, depois em Som.
  \item \textbf{Negrito}: muito importante! Pode ser um conceito-chave ou um aviso de segurança, preste muita atenção.
  \item \emph{Itálico}: algumas pequenas ideias complementares ou metáforas, para o ajudar a compreender melhor.
\end{itemize}

\booksection{III. Android e iOS, falamos de ambos}

Adivinhe, algumas coisas têm nomes diferentes, mas fazem o mesmo. Por exemplo: Aplicações (Android) e Apps (iOS), Definições (ambas se chamam Definições), Gestão de ficheiros (Android) e Ficheiros (iOS). Basta saber que é a mesma coisa, não se confunda com os termos.

Além disso, atualmente no campo Android existe também o sistema HarmonyOS da Huawei, com as suas próprias características, que abordaremos mais adiante. Em resumo, qualquer que seja o celular que utilize, este livro pode ajudá-lo.

\booksection{IV. Capacidade do celular: porque é que comprei 128G mas só tenho menos de 110G disponíveis?}

Todos conhecem as unidades de capacidade: 1TB = 1024GB, 1GB = 1024MB, 1MB = 1024KB... mas há uma armadilha: os fabricantes de discos costumam calcular em base 1000, enquanto o sistema calcula em base 1024, por isso o disco “128G” que comprou terá uma capacidade real “reduzida”.

Mas a redução da capacidade do seu celular deve-se principalmente a três razões:

\begin{enumerate}
  \item \textbf{O próprio sistema ocupa espaço}: o sistema operativo é como as paredes estruturais de uma casa, tem de existir. Hoje em dia, os sistemas costumam ocupar 10-20GB, o que é normal.
  \item \textbf{Aplicações pré-instaladas pelo fabricante}: as Apps que vêm no celular novo também ocupam espaço, algumas podem ser desinstaladas, outras não.
  \item \textbf{Perdas de partição e formatação}: por motivos técnicos, como o forro de uma mochila que também ocupa volume.
\end{enumerate}

\begin{tcbinfo}{Pequena sugestão de compra}
Hoje em dia, um celular de 128GB já não chega, 256GB é um ponto de partida mais seguro. Se gosta de fotografar ou jogar, recomendamos 512GB ou mais. Lembre-se, quanto ao espaço de armazenamento, mais vale ter a mais do que ficar aquém.

\end{tcbinfo}

\begin{tcbpractice}{Mãos à obra: veja a "identidade" do seu celular}
Abra [Definições], encontre Sobre o celular (ou opção semelhante) e anote: modelo do celular, versão do sistema, espaço de armazenamento total e disponível. Este é o primeiro passo para o conhecer. Aproveite para verificar o estado de saúde da bateria, para ter uma ideia.

\end{tcbpractice}

\booksection{V. O seu celular é uma “biblioteca de ativos digitais”}

Antes de começarmos os conhecimentos específicos, vamos estabelecer uma compreensão central:\textbf{o celular deixou há muito de ser apenas uma ferramenta de comunicação, é a nossa mais importante biblioteca de ativos digitais}. Contactos, histórico de conversas, fotografias de família, documentos de trabalho, Apps bancárias, códigos de pagamento, dados de saúde, documentos eletrónicos... todos estes dados juntos costumam valer mais do que o próprio celular.

\begin{tcbinfo}{Porque é que isto é importante?}
Se o celular se perder ou avariar, não perde apenas alguns milhares de yuans em hardware,\textbf{mas, mais ainda, dados irrecuperáveis acumulados durante anos}— fotografias do crescimento dos filhos, mensagens de voz de entes queridos falecidos, contratos importantes sem backup. Hábitos científicos de organização e proteção podem protegê-lo em três dimensões:

\begin{itemize}
  \item \textbf{Aumentar a eficiência}: uma classificação de ficheiros padronizada torna a pesquisa do sistema mais rápida e o celular não fica lento a longo prazo.
  \item \textbf{Prolongar a vida útil}: uma gestão de armazenamento adequada reduz a carga de leitura/escrita do chip, e tanto a bateria como o armazenamento duram mais.
  \item \textbf{Garantir a segurança dos dados}: gestão em camadas + backups regulares + restrição de permissões, mantém as "catástrofes de dados" afastadas.
\end{itemize}

\end{tcbinfo}

\booksubsection{Nossa filosofia central: “Prevenção em primeiro lugar, combinada com tratamento”}

Assim como cuidar da saúde não significa esperar ficar doente para consultar um médico,\textbf{a gestão do celular também deve começar pelos “hábitos diários”}. Este livro enfatiza repetidamente não uma técnica avançada específica, mas estes princípios simples:

\begin{itemize}
  \item \textbf{Prevenção em primeiro lugar}: Instale aplicativos apenas pelas lojas oficiais, remova metadados das fotos antes de enviar, não repita senhas — essas "pequenas ações" têm custo mínimo, mas bloqueiam a grande maioria dos riscos.
  \item \textbf{Check-up regular}: Organize a pasta de downloads uma vez por mês, faça uma auditoria de permissões a cada trimestre, um backup completo a cada semestre — um ritmo fixo é cem vezes mais eficaz do que apagar incêndios de improviso.
  \item \textbf{Ter um plano de contingência}: Ative a sincronização na nuvem, anote bem as contas e senhas, escreva uma "lista de herança digital" — garanta que você e sua família possam lidar com imprevistos com tranquilidade.
\end{itemize}
